 \chapter{Parcours et mobilités}\label{ch:in:career-paths-and-moves}

Un trader qui démissionne en début de février, quelques jours avant le versement de la prime de l'année précédente, peut devoir effectuer trois
mois de \omterm{def:in:career-paths-and-moves:notice}{préavis}, passer trois mois de plus en garden leave, puis respecter une clause de non-concurrence de six mois, et perdre en chemin la
prime de l'année précédente, la prime en cours d'acquisition pour l'année en cours et la rémunération différée non encore acquise. Dans le package illustratif de
ce chapitre, cela représente 659\,167~\$ de primes perdues et une année hors du marché. Le rachat offert par un nouvel employeur fait donc partie
de l'offre, et le moment de la démission vaut autant qu'une forte hausse salariale. Ce chapitre décrit comment les gens entrent dans
l'industrie et s'y déplacer, chiffre le coût d'un changement de poste, et suit une carrière comme une chaîne d'états.

\section{Voies d'entrée~: stages, programmes jeunes diplômés, recrutements latéraux, doctorats}

\begin{definition}[Stage, offre de retour, programme jeunes diplômés, recrutement latéral]\label{def:in:career-paths-and-moves:entry}
Un \emph{stage}\index{stage} est une période de travail rémunérée, de quelques semaines ou de quelques mois, destinée à un étudiant et conçue par
l'entreprise comme un essai préalable à un emploi permanent. Une \emph{offre de retour}\index{offre de retour} est l'offre d'un emploi permanent faite à un
stagiaire à la fin du stage. Un \emph{programme jeunes diplômés}\index{programme jeunes diplômés} est un parcours d'entrée structuré
proposé par l'entreprise aux nouveaux diplômés, avec une formation et souvent des rotations avant une affectation permanente. Un \emph{recrutement latéral}\index{recrutement latéral}
est une personne expérimentée recrutée auprès d'un autre employeur pour un poste de son niveau.
\end{definition}

Pour les diplômés, le \omterm{def:in:career-paths-and-moves:entry}{stage} est la principale porte d'entrée. Les entreprises le présentent ainsi~: un teneur de marché affirme que «~Notre
programme de \omterm{def:in:career-paths-and-moves:entry}{stages} est conçu pour préparer à un poste à temps plein~»~; un autre propose des \omterm{def:in:career-paths-and-moves:entry}{stages} en trading quantitatif,
en recherche, en apprentissage automatique et en génie logiciel dans ses bureaux de New York, Londres et Hong Kong. Le processus
de recrutement menant à un \omterm{def:in:career-paths-and-moves:entry}{stage} ou à un emploi de jeune diplômé (présélection, évaluations, entretiens et décision) fait l'objet
du livre~18.

\begin{dated}{2026-09}{Comment deux entreprises présentent leurs voies d'entrée}\label{dat:in:career-paths-and-moves:entry}
Optiver~: «~Notre programme de \omterm{def:in:career-paths-and-moves:entry}{stages} est conçu pour préparer à un poste à temps plein~»~; «~Nos postes de jeunes diplômés sont conçus autour
de l'apprentissage par le travail réel.~» Jane Street~: \omterm{def:in:career-paths-and-moves:entry}{stages} en «~trading quantitatif, apprentissage automatique, génie logiciel,
recherche quantitative, stratégie et produit~», dans ses bureaux de New York, Londres et Hong Kong.
\end{dated}

Le doctorat est une voie d'entrée à part entière pour les postes de recherche (chapitre~17)~: le chercheur rejoint l'entreprise après le doctorat,
parfois après un \omterm{def:in:career-paths-and-moves:entry}{stage} de recherche effectué pendant celui-ci, à un niveau supérieur à celui d'un jeune diplômé. Les recrutements latéraux couvrent le reste~: un
trader expérimenté, chercheur ou ingénieur qui passe d'une entreprise à l'autre, et dont ce chapitre chiffre le changement. Le chapitre~29 suit les
filières issues de la formation.

\section{De la banque au buy side et retour}

Les mobilités peuvent suivre n'importe quel edge de la carte de l'industrie. Des banques vers le buy side~: un \omterm{def:in:bank-quant:strat}{stratégiste de desk} ou un \omterm{def:in:structurer-sales-and-sales-trader:structurer}{structureur}
(chapitres~18 et~23) vers un hedge fund, un trader vers une plateforme (chapitre~22). Du buy side vers les banques, par exemple quand une
plateforme ferme un pod ou qu'un fonds ferme ses portes. Entre sociétés du buy side, un chercheur ou une équipe qui change d'employeur avec une stratégie, ce qui est
là où les clauses de non-concurrence et le droit du secret des affaires pèsent le plus (livre~16, chapitre~11). Entre entreprises technologiques et trading
(chapitre~20), dans les deux sens. Aucune donnée publique ne recense ces mouvements~; ce qui peut être mesuré, c'est leur coût.

\begin{definition}[Préavis]\label{def:in:career-paths-and-moves:notice}
Un \emph{préavis}\index{préavis} est le délai entre une démission ou un licenciement et la fin du contrat de travail,
fixé par le contrat, dans le respect du minimum légal~; le salarié est rémunéré et lié par le contrat pendant toute sa durée, et peut être
placé en garden leave pour tout ou partie de celle-ci.
\end{definition}

\begin{method}[Le coût d'un changement de poste]\label{met:in:career-paths-and-moves:move}
Soit un poste versant un \omterm{def:in:how-pay-works:base}{salaire de base} et une prime attendue $B$, versée au mois $p$ pour l'année précédente, dont une part $d$ est différée en $V$
tranches annuelles. En régime stationnaire, les différés non acquis valent $dB(V+1)/2$. Démissionner au mois $m$ entraîne leur perte,
la prime de l'année précédente si $m\le p$, et la prime en cours d'acquisition de l'année, soit $B(m-1)/12$. Hors du marché pendant $T$ mois
(\omterm{def:in:career-paths-and-moves:notice}{préavis}, garden leave, clause de non-concurrence), celui qui change de poste perd $BT/12$ de prime, ainsi que le \omterm{def:in:how-pay-works:base}{salaire de base} des mois non payés. Un rachat
remplace une part $s$ des primes perdues avec la probabilité $q$. Le changement est neutre lorsqu'une hausse $\rho$ sur le \omterm{def:in:how-pay-works:base}{salaire de base} et la prime
sur $H$ années compense le coût net $C$~: $\rho=C/((\text{base}+B)H)$.
\end{method}

\begin{example}[Changer maintenant ou attendre le versement~?]\label{ex:in:career-paths-and-moves:move}
Conditions illustratives~: \omterm{def:in:how-pay-works:base}{salaire de base} 250\,000~\$, prime attendue 350\,000~\$, 40~\% différés sur trois ans, primes versées en février~;
\omterm{def:in:career-paths-and-moves:notice}{préavis} de trois mois, trois mois de garden leave et une clause de non-concurrence de six mois rémunérée~; un rachat de toutes les primes perdues avec
une probabilité d'un demi~; un horizon de trois ans. Les différés non acquis valent 280\,000~\$. Démissionner en début de février
fait perdre 659\,167~\$ et coûte, net du rachat attendu, 694\,167~\$~: une hausse de 38.6~\% sur le package de 600\,000~\$. Attendre
un mois, jusqu'au versement de la prime, ramène la perte à 338\,333~\$ et la hausse neutre à 30.5~\%. Sans rachat
le changement de février exige 56.1~\%~; avec un rachat certain, 21.1~\% (\cref{fig:in:career-paths-and-moves:surface}).
\end{example}

\begin{omfigure}
\begin{tikzpicture}
\begin{axis}[width=10.5cm,height=5.6cm,xmin=0.5,xmax=12.5,ymin=0,ymax=60,xtick={1,2,3,4,5,6,7,8,9,10,11,12},
  xticklabels={janv.,févr.,mars,avr.,mai,juin,juil.,août,sept.,oct.,nov.,déc.},xticklabel style={font=\scriptsize},typeset ticklabels with strut,xlabel={\small mois de la démission},
  ylabel={\small hausse neutre, \%},tick label style={font=\footnotesize},grid=major,grid style={gray!20},
  table/col sep=comma,legend style={legend cell align=left,font=\scriptsize,at={(0.5,-0.3)},anchor=north,legend columns=5,
  /tikz/every even column/.append style={column sep=6pt}}]
\addplot[black!60,thick,mark=o,mark options={solid}] table[x=month,y=q0]{figdata/industry/28-career-paths-and-moves/surface.csv};
\addplot[omThm,thick,mark=square*] table[x=month,y=q25]{figdata/industry/28-career-paths-and-moves/surface.csv};
\addplot[omDef,thick,mark=*] table[x=month,y=q50]{figdata/industry/28-career-paths-and-moves/surface.csv};
\addplot[omDef!50,thick,mark=triangle*] table[x=month,y=q75]{figdata/industry/28-career-paths-and-moves/surface.csv};
\addplot[black,thick,dashed,mark=diamond*,mark options={solid}] table[x=month,y=q100]{figdata/industry/28-career-paths-and-moves/surface.csv};
\legend{sans rachat,$q=0.25$,$q=0.5$,$q=0.75$,rachat certain}
\end{axis}
\end{tikzpicture}

\omcaption{La hausse sur le \omterm{def:in:how-pay-works:base}{salaire de base} et la prime qui rend neutre le changement illustratif du chapitre sur trois ans, selon le mois de
démission et la probabilité $q$ d'un rachat total des primes perdues. La prime est versée en février~: une démission antérieure au
versement la fait perdre, et la prime en cours d'acquisition de l'année croît chaque mois suivant. Données~: \texttt{in\_career.surface}, à partir de
\texttt{firm.careerpath.breakeven\_raise}.}\label{fig:in:career-paths-and-moves:surface}
\end{omfigure}

La forme de la \cref{fig:in:career-paths-and-moves:surface} est la leçon pratique. Le coût chute brutalement dès qu'une prime est
versée, puis remonte au cours de l'année à mesure que la suivante s'accumule~; un rachat aplatit la chute au moment du versement, car il
remplace les primes perdues, mais ni la prime en cours d'acquisition ni le temps passé hors du marché. L'année hors du marché coûte 350\,000~\$ de primes
à chaque mois~: la clause de non-concurrence est le poste le plus important.

\section{Clauses de non-concurrence et garden leave en pratique}

Le livre~16, chapitre~11, expose le droit~: clauses de non-concurrence, garden leave, non-sollicitation et confidentialité, et la manière
dont les tribunaux les traitent. La pratique dépend du lieu d'exercice du poste.

\begin{dated}{2026-09}{Droit du préavis et des clauses de non-concurrence}\label{dat:in:career-paths-and-moves:law}
Royaume-Uni~: \omterm{def:in:career-paths-and-moves:notice}{préavis} légal minimal, de la part de l'employeur, d'une semaine par année d'ancienneté dans la limite de douze semaines, et de la part
du salarié, une semaine (Employment Rights Act 1996, s.~86)~; un contrat peut fixer des durées plus longues. Le document de travail
du gouvernement du 26 novembre 2025 a été soumis à consultation jusqu'au 18 février 2026, sur des options comprenant une limite légale de la durée
des clauses de non-concurrence, une interdiction, ou une interdiction en dessous d'un seuil de salaire~; le gouvernement précédent avait annoncé une limite de trois mois
le 10 mai 2023, qui n'a jamais été adoptée. États-Unis~: la règle de la FTC sur les clauses de non-concurrence a été bloquée par un tribunal de district le
20 août 2024, et le 5 septembre 2025 la Commission a voté le retrait de ses recours et accepté l'annulation~; la
règle n'est pas en vigueur. Californie~: une clause de non-concurrence nulle selon son droit «~n'est pas exécutoire, quels que soient le lieu et le moment où le
contrat a été signé~» (Business and Professions Code s.~16600.5).
\end{dated}

Le calcul côté employeur repose sur \texttt{firm.gardenleave} du livre~16~: payer quelqu'un pour qu'il reste hors du marché
vaut la peine tant que le savoir sur les stratégies qu'il emporte se déprécie moins vite que ne coûte sa rémunération. Avec les données illustratives du chapitre
(une stratégie rapportant 5~millions de dollars par an, 20~\% de ce gain perdu au profit d'un concurrent qui l'apprend, une demi-vie de l'edge d'un an, et 600\,000~\$ de salaire par an pendant le congé)
la meilleure durée est de 0.74 an, soit environ neuf mois, ce qui est l'ordre de grandeur
des congés que suppose l'exemple. Une entreprise qui paie un congé achète du temps pour faire évoluer ses stratégies.

\section{Se mettre à son compte et créer une société}

Certains mouvements sortent complètement de l'emploi~: un \omterm{def:in:portfolio-manager-and-pod-analyst:roles}{gérant de portefeuille} lance un fonds, un groupe de traders crée une société de trading pour compte propre, un ingénieur fonde un prestataire.
Le livre~16, chapitre~25, traite du lancement d'un fonds et de la levée de capitaux, avec les programmes pour gérants émergents et les investisseurs d'amorçage
qui soutiennent les nouveaux fonds, et de la portabilité d'un historique de performance, qui détermine si les investisseurs y croient.
Le chapitre~3, sur la généalogie des sociétés de trading, montre comment plusieurs teneurs de marché d'options ont débuté sur les parquets des bourses~;
il ne donne pas de données sourcées sur les \omterm{def:in:the-options-market-making-houses:spinout}{essaimages} (spin-out), car le registre public documente rarement qui a quitté quelle société pour en fonder une autre.
L'économie de ce mouvement est celle de l'associé du chapitre~25, sans le salaire~: un capital exposé dès le premier jour
et un revenu qui ne vient que des profits.

\section{Sorties}

Des carrières quittent aussi le secteur~: vers les entreprises technologiques, vers le monde académique, vers le secteur public
(régulateurs et banques centrales, chapitre~9), vers la trésorerie et le risque des entreprises, ou vers la retraite après une courte carrière.
La chaîne de carrière illustrative de la \cref{fig:in:career-paths-and-moves:fan} donne à chaque année 5~\% de chances de quitter le secteur~; avec ses autres hypothèses, 36.7~\%
des diplômés sont partis au bout de dix ans et 21.5~\% gèrent un portefeuille. Leurs trajectoires de salaire se dispersent fortement~: la médiane atteint le salaire d'un chercheur senior
dès la huitième année, puis redescend à mesure que les sorties s'accumulent, tandis que le 90\up{e}~centile rejoint celui
d'un \omterm{def:in:portfolio-manager-and-pod-analyst:roles}{gérant de portefeuille} dès la septième année.

\begin{omfigure}
\begin{tikzpicture}
\begin{axis}[width=10.5cm,height=5.6cm,xmin=1,xmax=15,ymin=0,ymax=1800,xlabel={\small année de carrière},
  ylabel={\small salaire annuel, en milliers de \$},tick label style={font=\footnotesize},grid=major,grid style={gray!20},
  table/col sep=comma,legend style={font=\scriptsize,at={(0.03,0.97)},anchor=north west,legend cell align=left}]
\addplot[name path=hi,draw=none,forget plot] table[x=year,y=p90]{figdata/industry/28-career-paths-and-moves/fan.csv};
\addplot[name path=lo,draw=none,forget plot] table[x=year,y=p10]{figdata/industry/28-career-paths-and-moves/fan.csv};
\addplot[omDef!20] fill between[of=hi and lo];
\addplot[omDef,thick] table[x=year,y=mean]{figdata/industry/28-career-paths-and-moves/fan.csv};
\addplot[omThm,thick,dashed] table[x=year,y=p50]{figdata/industry/28-career-paths-and-moves/fan.csv};
\legend{du 10\up{e} au 90\up{e}~centile,moyenne,médiane}
\end{axis}
\end{tikzpicture}

\omcaption{Salaire annuel le long d'une chaîne de carrière illustrative (diplômé, chercheur, chercheur senior, \omterm{def:in:portfolio-manager-and-pod-analyst:roles}{gérant de portefeuille},
sorti du secteur) avec des probabilités de transition étiquetées et le salaire par état~: moyenne, médiane et bande comprise entre le 10\up{e} et le 90\up{e}~centile
sur 20\,000 carrières. Chaque chiffre est une hypothèse~; la forme, un éventail qui s'élargit, est l'essentiel. Données~:
\texttt{in\_career.career}, tirées de \texttt{firm.careerpath.simulate}.}\label{fig:in:career-paths-and-moves:fan}
\end{omfigure}

\section{Tutoriel~: partir maintenant ou attendre l'acquisition des droits~?}\label{tut:in:career-paths-and-moves:tut}

\begin{tutorial}
\textbf{Objectif.} Évaluer un départ selon son mois et son rachat, et suivre une carrière comme une chaîne de Markov.
\textbf{État final~:} \cref{fig:in:career-paths-and-moves:surface,fig:in:career-paths-and-moves:fan}.
\begin{tutsteps}
\item \textbf{Le package.} \texttt{firm.careerpath.Job(base, award, deferral, vest\_years, pay\_month)}~; \texttt{forfeited}
renvoie, par mois, les trois éléments frappés de \omterm{def:in:how-pay-works:rsu}{déchéance}.
\item \textbf{Le départ.} \texttt{move\_cost} et \texttt{breakeven\_raise} (\cref{lst:in:career-paths-and-moves:move}).
\omcode{code/firm/careerpath/firm_careerpath.py}{50}{62}{Le coût net d'un départ et la hausse qui le rend neutre.}\label{lst:in:career-paths-and-moves:move}
\item \textbf{Le point de vue de l'employeur.} \texttt{firm.gardenleave.best\_length} du livre~16.
\item \textbf{La carrière.} \texttt{Chain(states, P, pay)}, \texttt{simulate} et \texttt{pay\_paths}~; les \omterm{def:in:the-options-market-making-houses:spinout}{essaimages} du graphe de
généalogie, via \texttt{spinouts}.
\end{tutsteps}
\end{tutorial}

Les points extrêmes de la surface~: démissionner en début de février exige 56.1~\% sans rachat et 21.1~\% avec un rachat certain~;
démissionner en début de décembre, 52.8~\% et 37.3~\%.

\textbf{Pour aller plus loin.} Actualiser la hausse et laisser varier l'attribution~; faire couvrir au seul rachat le salaire différé~; ajouter
une clause de good-leaver (chapitre~13)~; estimer les probabilités de transition d'une chaîne à partir de profils publics, avec leur biais.

\section{Construction~: firm.careerpath}\label{bld:in:career-paths-and-moves:careerpath}

\begin{build}
\bfield{Rôle} Chiffrer un départ entre employeurs, et modéliser une carrière comme une chaîne d'états dont les hypothèses sont étiquetées.
\bfield{Interface} \texttt{firm.careerpath}~: \texttt{Job}~; \texttt{unvested}~; \texttt{forfeited(job, month)}~;
\texttt{move\_cost(job, month, notice\_m, garden\_m, noncompete\_m, paid\_noncompete, q, s)}~;
\texttt{breakeven\_raise(job, month, horizon\_y, \dots)}~; \texttt{Chain(states, P, pay)}~;
\texttt{simulate(chain, start, years, n, rng)}~; \texttt{pay\_paths}~; \texttt{spinouts(graph)}~; avec \texttt{firm.gardenleave} et \texttt{firm.lineage}.
\bfield{Règles} \omterm{def:in:how-pay-works:rsu}{Déchéance} pour bad-leaver~; aucune attribution ne s'accumule hors du marché~; les lignes de la chaîne somment à un, ou la simulation
refuse~; chaque paramètre est étiqueté par l'appelant.
\bfield{Tests d'acceptation} \texttt{code/firm/careerpath/tests/}~: l'attribution non versée n'est frappée de \omterm{def:in:how-pay-works:rsu}{déchéance} qu'avant son
mois de versement~; le pire mois est celui qui précède le versement~; un rachat certain compense les attributions frappées de \omterm{def:in:how-pay-works:rsu}{déchéance}~; une chaîne déterministe
évolue comme spécifié.
\bfield{Pour aller plus loin} Actualisation et incertitude sur l'attribution (avec \texttt{firm.payoffer})~; une chaîne dont la
promotion dépend de la durée~; mouvements d'équipes.
\end{build}

\begin{omsources}
\item Employment Rights Act 1996, s.~86~; UK Department for Business and Trade (2025), document de travail sur la réforme des clauses de non-concurrence.
\item Federal Trade Commission (2025), communiqué de presse sur la règle relative aux clauses de non-concurrence~; California Business and Professions Code
s.~16600.5.
\item Optiver et Jane Street, pages «~students~» et «~internships~».
\item le livre~16, chapitres~11 et~25~; le livre~18 sur le recrutement.
\end{omsources}

\section{Exercices}

\begin{exercise}[$\star$]\label{exo:in:career-paths-and-moves:1}
Calculer, avec le package de l'exemple, les montants différés non acquis, puis déterminer ce qu'une démission en début de mars fait perdre par \omterm{def:in:how-pay-works:rsu}{déchéance}.
\end{exercise}

\begin{exercise}[$\star$]\label{exo:in:career-paths-and-moves:2}
Calculer le \omterm{def:in:career-paths-and-moves:notice}{préavis} légal minimal au Royaume-Uni qu'un employeur doit donner à un salarié ayant sept ans d'ancienneté.
\end{exercise}

\begin{exercise}[$\star$]\label{exo:in:career-paths-and-moves:3}
Déterminer le montant d'attribution que coûte une année hors du marché dans l'exemple.
\end{exercise}

\begin{exercise}[$\star\star$]\label{exo:in:career-paths-and-moves:4}
Un rachat certain rend un départ en février moins coûteux qu'un départ en mars sans rachat. Expliquer pourquoi il n'est pas pour autant gratuit.
\end{exercise}

\begin{exercise}[$\star\star$]\label{exo:in:career-paths-and-moves:5}
Un nouvel employeur propose une hausse de 25~\% et une probabilité de rachat de 50~\%. Déterminer dans quels mois le départ en vaut la peine.
\end{exercise}

\begin{exercise}[$\star\star$]\label{exo:in:career-paths-and-moves:6}
Expliquer pourquoi un candidat aurait intérêt à savoir si une clause de non-concurrence est rémunérée.
\end{exercise}

\begin{exercise}[$\star\star\star$]\label{exo:in:career-paths-and-moves:7}
\emph{Programmation.} Rendre la clause de non-concurrence non rémunérée. Calculer la hausse neutre pour une démission en février avec
une probabilité de rachat de 50~\%.
\end{exercise}

\begin{exercise}[$\star\star\star$]\label{exo:in:career-paths-and-moves:8}
\emph{Trouver la faille.} «~Le modèle montre qu'un cinquième des diplômés deviennent \omterm{def:in:portfolio-manager-and-pod-analyst:roles}{gérants de portefeuille} en dix ans.~»
\end{exercise}

\section{Problème~: partir maintenant ou attendre l'acquisition des droits}

\begin{problem}[{Problème du week-end --- partir maintenant ou attendre l'acquisition des droits~?}]\label{pb:in:career-paths-and-moves:1}
Un trader, dont le package est celui de l'exemple, reçoit une offre fin janvier et demande à quel moment démissionner.

\medskip\noindent\textbf{Partie I --- Les parcours.}
\begin{enumerate}
\item Définir un \omterm{def:in:career-paths-and-moves:entry}{stage}, une \omterm{def:in:career-paths-and-moves:entry}{offre de retour}, un \omterm{def:in:career-paths-and-moves:entry}{programme jeunes diplômés} et un \omterm{def:in:career-paths-and-moves:entry}{recrutement latéral}.
\item Décrire la manière dont les sociétés présentent leurs \omterm{def:in:career-paths-and-moves:entry}{stages} et leurs postes de jeunes diplômés.
\item Indiquer où mène un doctorat.
\item Nommer les edges communs des mobilités entre employeurs.
\item Indiquer quelles données publiques les recensent.
\end{enumerate}

\medskip\noindent\textbf{Partie II --- Le coût.}
\begin{enumerate}[resume]
\item Définir le \omterm{def:in:career-paths-and-moves:notice}{préavis} et énoncer le minimum légal au Royaume-Uni.
\item Énoncer la méthode.
\item Calculer la \omterm{def:in:how-pay-works:rsu}{déchéance} et le coût net pour une démission en février et une démission en mars.
\item Calculer la hausse neutre dans chacun des cas, avec une probabilité de rachat de 50~\%.
\item Identifier l'élément le plus important pris isolément, et expliquer pourquoi.
\end{enumerate}

\medskip\noindent\textbf{Partie III --- Le droit et l'employeur.}
\begin{enumerate}[resume]
\item Indiquer le statut de la règle de la FTC sur la non-concurrence.
\item Indiquer ce que le gouvernement britannique a proposé, et quand.
\item Indiquer ce qu'ajoute la section s.~16600.5 du droit de Californie.
\item Déterminer la durée de garden leave que donne le modèle de l'employeur, et à partir de quand.
\item Indiquer ce dont un \omterm{def:in:the-options-market-making-houses:spinout}{essaimage} a besoin et qu'un simple départ n'exige pas.
\end{enumerate}

\medskip\noindent\textbf{Partie IV --- Le verdict.}
\begin{enumerate}[resume]
\item Énoncer le \emph{résultat nommé}~: la hausse qui rend le départ neutre en fonction du mois de départ et de la
probabilité de rachat, et sa valeur la veille du versement de l'attribution.
\item Indiquer ce qu'il doit demander au nouvel employeur de racheter.
\item Expliquer comment une \omterm{def:in:how-pay-works:rsu}{clause de good leaver} modifierait la réponse.
\item Indiquer ce que dit la chaîne de carrière, et ce qu'elle ne dit pas.
\item En deux phrases, conseiller-le.
\end{enumerate}
\end{problem}

\section{Questions d'entretien}

\begin{interviewq}[$\star$ \iqroles{trader, chercheur}]\label{iq:in:career-paths-and-moves:1}
Pourquoi souhaitez-vous quitter votre société actuelle~?
\end{interviewq}

\begin{interviewq}[$\star$ \iqroles{chercheur}]\label{iq:in:career-paths-and-moves:2}
Que pouvez-vous nous dire de vos stratégies actuelles, et que ne pouvez-vous pas dire~?
\end{interviewq}

\begin{interviewq}[$\star\star$ \iqroles{trader}]\label{iq:in:career-paths-and-moves:3}
Vous avez une clause de non-concurrence de six mois. Comment utiliseriez-vous ce temps~?
\end{interviewq}

\begin{interviewq}[$\star\star$ \iqroles{chercheur, trader}]\label{iq:in:career-paths-and-moves:4}
Comment nous montreriez-vous votre historique de performance sans divulguer les informations de votre employeur~?
\end{interviewq}

\begin{interviewq}[$\star\star$ \iqroles{banque, chercheur}]\label{iq:in:career-paths-and-moves:5}
Pourquoi quitter une banque pour un fonds, et que regretterez-vous~?
\end{interviewq}

\begin{interviewq}[$\star\star\star$ \iqroles{trader, chercheur}]\label{iq:in:career-paths-and-moves:6}
Évaluez la rémunération différée que vous abandonneriez en nous rejoignant, et proposez un rachat.
\end{interviewq}
